\section{Summary}

The study asks how instance characteristics relate to the performance of
algorithms for the static H-ODBRP. Instances are described by demand, time
window, fleet, and network features, then projected into a low-dimensional
instance space. SA and GR form the present performance matrix. The exact
Gurobi model is reserved for validation on very small cases; a MIP-based
large-neighbourhood matheuristic is intended to show where SA can improve, not
to act as a third ISA algorithm.

Work completed includes a 72-case Copenhagen pilot, a correction to request
seeding, a network cut-out for Bornholm, a nearest-node cache, solver input
fixes, an exact-model checker, a rewritten exact model, and HPC job scripts.
The pilot suggests that SA is usually better than GR: SA wins on 70 of 72
cases. GR is competitive in a narrow corner: tight time windows and the
smallest pilot fleet. The pilot projection retains a gap associated with the
two time-window settings. The available feature CSVs do not independently
reproduce the documented pre-fix count of six distinct mean OD distances; the
post-fix CSV does contain 24.

The proposed full design covers 100--1000 requests in a two-hour period and
draws instance settings and fleets using Latin hypercubes. The region, period,
and ranges remain \tagprop; the author has the final decision. The new exact
model agrees with the heuristic feasibility checks in all 11 tested
scenarios, while the original model disagreed in 9. On rerun, the new model
proved the five-request case in 2.5 seconds, but had a 25\% gap at the limit
for 10 requests and did not prove the 15-request case.

The full instance set is still blocked by unresolved generation cost and
memory questions, an unrun station-matrix build, and the missing shared-matrix
feature step. Other open work includes the third ISA algorithm, the
matheuristic HPC test, solver resource sizing, and final design choices. These
results are preliminary, and the one-sided SA--GR result must be tested on the
wide design.
